% Created 2024-11-21 Thu 19:36
% Intended LaTeX compiler: pdflatex
\documentclass[11pt]{article}
\usepackage[utf8]{inputenc}
\usepackage[T1]{fontenc}
\usepackage{graphicx}
\usepackage{longtable}
\usepackage{wrapfig}
\usepackage{rotating}
\usepackage[normalem]{ulem}
\usepackage{amsmath}
\usepackage{amssymb}
\usepackage{capt-of}
\usepackage{hyperref}
\author{Prajwal Rao}
\date{\today}
\title{Spearman correlation}
\hypersetup{
 pdfauthor={Prajwal Rao},
 pdftitle={Spearman correlation},
 pdfkeywords={},
 pdfsubject={},
 pdfcreator={Emacs 29.1 (Org mode 9.7.11)}, 
 pdflang={English}}
\usepackage{biblatex}

\begin{document}

\maketitle
\tableofcontents

Spearman correlation measures the strength and direction of association between two ranked variables. It is a non-parametric measure, which means it does not assume a linear relationship or require the data to be normally distributed. Instead, it assesses how well the relationship between two variables can be described by a monotonic function.
\subsubsection{Mathematical Explanation:}
\label{sec:org4cfec65}

\begin{enumerate}
\item \textbf{Rank the Data}: Convert the data values of each variable to ranks.
\item \textbf{Difference of Ranks}: For each pair of ranks, compute the difference \(d_i\) between the two ranks.
\item \textbf{Square the Differences}: Compute \(d_i^2\) for each pair of ranks.
\item \textbf{Sum of Squares}: Calculate the sum \(\sum d_i^2\).
\item \textbf{Spearman's Correlation Coefficient} \(\rho\) is computed using the formula:

\[
   \rho = 1 - \frac{6 \sum d_i^2}{n(n^2 - 1)}
   \]

where \(n\) is the number of paired ranks.
\end{enumerate}
\subsubsection{Python Example:}
\label{sec:org37ff683}

Here’s how to compute Spearman correlation using Python with the \texttt{scipy} library:

\begin{verbatim}
import numpy as np
from scipy.stats import spearmanr

# Sample data
x = [1, 2, 3, 4, 5]
y = [5, 6, 7, 8, 7]

# Compute Spearman correlation
correlation, p_value = spearmanr(x, y)

print(f"Spearman correlation: {correlation}, p-value: {p_value}")
\end{verbatim}
\subsubsection{Output:}
\label{sec:orgc14d131}
This script will output the Spearman correlation coefficient and the p-value, indicating the strength and statistical significance of the correlation between the two variables.
\section{For Images}
\label{sec:org8f54237}
Spearman correlation assesses the rank-order relationship between two variables. For images, this involves treating the pixel values as one-dimensional arrays and computing the Spearman correlation between them.
\subsubsection{Steps:}
\label{sec:org38e5bef}

\begin{enumerate}
\item \textbf{Flatten the Images:} Convert the images into 1D arrays.
\item \textbf{Rank the Values:} Assign ranks to the pixel values in both arrays.
\item \textbf{Calculate Spearman Correlation:} Use the formula:
\[
   \rho = 1 - \frac{6 \sum d_i^2}{n(n^2 - 1)}
   \]
where \(d_i\) is the difference between the ranks of corresponding pixels, and \(n\) is the number of pixels.
\end{enumerate}
\subsubsection{Python Example:}
\label{sec:org31e6003}

You can use libraries like \texttt{numpy} and \texttt{scipy} to calculate the Spearman correlation.

\begin{verbatim}
import numpy as np
from scipy.stats import spearmanr
from PIL import Image

# Load two images
image1 = Image.open('image1.png').convert('L')  # Convert to grayscale
image2 = Image.open('image2.png').convert('L')  # Convert to grayscale

# Flatten the images to 1D arrays
data1 = np.array(image1).flatten()
data2 = np.array(image2).flatten()

# Calculate Spearman correlation
correlation, _ = spearmanr(data1, data2)

print(f"Spearman correlation: {correlation}")
\end{verbatim}
\subsubsection{Explanation of Code:}
\label{sec:orgbdab6b0}
\begin{itemize}
\item \textbf{Image Loading:} The images are loaded and converted to grayscale.
\item \textbf{Flattening:} The pixel values are turned into 1D arrays using \texttt{.flatten()}.
\item \textbf{Spearman Calculation:} The \texttt{spearmanr()} function from \texttt{scipy.stats} computes the correlation.
\end{itemize}

This method can effectively measure the monotonic relationship between the pixel intensity distributions of the two images.
\end{document}
